\section{Task 1 -- Khai thác và đặc tả yêu cầu}\label{ch:yeucau}

Tạo phẩm của Task 1 nằm trong \code{docs/\allowbreak{}requirements/\allowbreak{}}:
\begin{itemize}
\item \code{context-and-scope.\allowbreak{}md}: bối cảnh, stakeholder, phạm vi.
\item \code{srs.md}: đặc tả yêu cầu.
\item \code{use-case-checkout.\allowbreak{}md}: đặc tả UC-01.
\item Hai sơ đồ use case dạng PlantUML kèm ảnh kết xuất.
\end{itemize}

Task 1 được merge qua PR \#2, \#3 và \#4 ngày 14/11/2026, trước mọi tạo phẩm của Task 2.

\subsection{Task 1.1 -- Bối cảnh, stakeholder và phạm vi}

\subsubsection{Vấn đề nghiệp vụ}
Quán cà phê cần một kênh đặt món và thanh toán không tiền mặt để giảm xếp hàng và nhầm đơn. Có ba điểm nghiệp vụ khó, quyết định thiết kế của toàn hệ thống:

\begin{enumerate}
\item \textbf{Không bán vượt nguyên liệu.} Mỗi món được pha từ nhiều nguyên liệu theo công thức, và lượng dùng thay đổi theo size, mức đường, mức đá, topping.
\item \textbf{Không trừ tiền hai lần} khi khách bấm thanh toán lại hoặc mạng chập chờn.
\item \textbf{Đơn hàng đi đúng vòng đời}, từ lúc đặt đến lúc giao.
\end{enumerate}

\subsubsection{Stakeholder}
Đề bài nêu sáu nhóm stakeholder. Bảng \ref{tab:stakeholder} ghi mục tiêu của từng nhóm và yêu cầu chịu ảnh hưởng.

{\small
\begin{longtable}{@{}L{2.6cm}L{2.4cm}L{4.6cm}L{4.3cm}@{}}
\caption{Stakeholder, mục tiêu và ảnh hưởng đến yêu cầu}\label{tab:stakeholder}\\
\toprule
\textbf{Stakeholder} & \textbf{Vai trò} & \textbf{Mục tiêu và kỳ vọng} & \textbf{Ảnh hưởng đến yêu cầu} \\ \midrule \endfirsthead
\toprule \textbf{Stakeholder} & \textbf{Vai trò} & \textbf{Mục tiêu và kỳ vọng} & \textbf{Ảnh hưởng đến yêu cầu} \\ \midrule \endhead
\bottomrule \endfoot
Khách hàng & Đặt và thanh toán đơn & Đặt nhanh, giá rõ ràng, biết món còn hay hết, theo dõi được đơn, được thưởng khi mua nhiều & Menu, giỏ hàng, thanh toán, lịch sử đơn, điểm thưởng, đổi quà \\
Nhân viên quầy & Pha chế, giao món & Chỉ nhận đơn đã thanh toán, cập nhật trạng thái nhanh, biết nguyên liệu sắp hết & Màn hình xử lý đơn, máy trạng thái đơn, quản lý nguyên liệu \\
Quản lý cửa hàng & Kiểm soát vận hành & Không bán vượt nguyên liệu đang có, kiểm soát khuyến mãi và điểm thưởng & Công thức và tồn kho nguyên liệu, khuyến mãi, quà đổi điểm \\
Đơn vị thanh toán (cổng mô phỏng, VNPay sandbox) & Xử lý Ví, Thẻ, QR & Nhận yêu cầu hợp lệ, không trừ tiền hai lần & Idempotency, trạng thái thanh toán, hạn của mã QR \\
Giảng viên (khách hàng nghiệm thu) & Chấm và nghiệm thu & Sản phẩm chạy được, đủ tạo phẩm của Task 1 đến Task 5, quy trình đúng thứ tự & Toàn bộ tạo phẩm Assignment 2 \\
Nhóm phát triển & Xây dựng sản phẩm & Xong đúng hạn, mã dễ bảo trì, có CI/CD & Ràng buộc công nghệ, quy trình Scrum, CI/CD \\
\end{longtable}
}

\subsubsection{Phạm vi}
Phạm vi được chia theo MoSCoW (bảng \ref{tab:phamvi} ở mục \ref{ch:gioithieu}).
\begin{itemize}
\item \textbf{MVP:} chỉ gồm các chức năng cần cho luồng nghiệp vụ tối thiểu của đề.
\item \textbf{Chức năng bổ sung:} mã khuyến mãi, điểm thưởng, kho nguyên liệu theo công thức, thanh toán QR, đổi quà, tự hủy đơn, VNPay sandbox. Các chức năng này phục vụ yêu cầu của giảng viên là sản phẩm phải giống một trang bán hàng thật, có logic hoàn chỉnh. Chúng được làm sau MVP, theo thứ tự ưu tiên.
\item \textbf{Không làm (Won't):} được ghi rõ để tránh mở rộng phạm vi trong lúc làm.
\end{itemize}

\subsubsection{Giả định, ràng buộc và rủi ro}
Những điểm đề bài chưa nêu rõ được nhóm chốt thành giả định (bảng \ref{tab:giadinh}). Quan trọng nhất là các giả định về giá, hệ số nguyên liệu theo size, điểm thưởng và thời điểm trừ kho.

{\small
\begin{longtable}{@{}L{1cm}L{12.9cm}@{}}
\caption{Giả định}\label{tab:giadinh}\\
\toprule
\textbf{\#} & \textbf{Giả định} \\ \midrule \endfirsthead
\toprule \textbf{\#} & \textbf{Giả định} \\ \midrule \endhead
\bottomrule \endfoot
G1 & Thanh toán là mô phỏng. Kết quả thành công hay thất bại điều khiển được để kiểm thử. Tài khoản thanh toán thử nghiệm và số dư là dữ liệu seed, không phải tiền thật \\
G2 & Một cửa hàng, một loại tiền (VND) \\
G3 & Phụ phí size: S +0, M +5.000, L +10.000 đồng \\
G4 & Lượng nguyên liệu theo size: S 100\%, M 125\%, L 150\% công thức, làm tròn lên \\
G5 & Điểm thưởng: 1 điểm cho mỗi 10.000 đồng của đơn đã thanh toán; hủy đơn đã thanh toán thì trừ lại điểm \\
G6 & Tạo đơn chỉ kiểm tra đủ nguyên liệu. Nguyên liệu bị trừ khi thanh toán thành công. Hủy đơn chỉ hoàn phần đã trừ \\
G7 & Đơn mới ở trạng thái PENDING; thanh toán chuyển đơn sang PAID hoặc PAYMENT\_FAILED \\
G8 & STAFF gồm cả nhân viên và quản lý; không có vai trò ADMIN riêng \\
G9 & Với các điểm đề bài chưa nêu rõ, nhóm tự chọn cách làm và báo giảng viên xác nhận \\
\end{longtable}
}
{\small
\begin{longtable}{@{}L{1cm}L{12.9cm}@{}}
\caption{Ràng buộc}\label{tab:rangbuoc}\\
\toprule
\textbf{\#} & \textbf{Ràng buộc} \\ \midrule \endfirsthead
\toprule \textbf{\#} & \textbf{Ràng buộc} \\ \midrule \endhead
\bottomrule \endfoot
R1 & Bắt buộc dùng Next.js (frontend), NestJS (backend) và PostgreSQL; được bổ sung công nghệ khác \\
R2 & Frontend theo hướng mobile-first \\
R3 & Hạn nộp: 15/11/2026 \\
R4 & Nhóm 4 người, làm theo Scrum, chia hai sprint \\
R5 & Repository Assignment 2 độc lập hoàn toàn với Assignment 1: không dùng chung mã nguồn, tài liệu hay lịch sử commit \\
R6 & Các Task làm tuần tự. Task sau chỉ bắt đầu khi tạo phẩm của Task trước đã có trên repository \\
R7 & Production truy cập qua HTTPS; database production không public; secret không nằm trong repository \\
\end{longtable}
}
{\small
\begin{longtable}{@{}L{1cm}L{4.2cm}L{1.7cm}L{1.7cm}L{5.3cm}@{}}
\caption{Rủi ro ban đầu và cách giảm thiểu}\label{tab:ruiro}\\
\toprule
\textbf{\#} & \textbf{Rủi ro} & \textbf{Khả năng} & \textbf{Ảnh hưởng} & \textbf{Giảm thiểu} \\ \midrule \endfirsthead
\toprule \textbf{\#} & \textbf{Rủi ro} & \textbf{Khả năng} & \textbf{Ảnh hưởng} & \textbf{Giảm thiểu} \\ \midrule \endhead
\bottomrule \endfoot
RK1 & Thời gian còn rất ngắn, không kịp đủ chức năng bổ sung & Cao & Cao & Làm xong MVP (Must) trước; chức năng Should và Could làm sau, theo thứ tự ưu tiên \\
RK2 & Không kịp triển khai production & Trung bình & Cao & Dựng pipeline CD và deploy thử ngay khi có MVP đầu tiên \\
RK3 & Đặt hoặc thanh toán đồng thời làm âm tồn kho, trừ tiền hai lần & Trung bình & Cao & Cập nhật có điều kiện trong transaction, Idempotency-Key, có test đồng thời \\
RK4 & Sơ đồ, SRS và mã nguồn mâu thuẫn nhau & Trung bình & Trung bình & Duy trì \code{docs/\allowbreak{}traceability-matrix.\allowbreak{}md} từ Task 2 \\
RK5 & Thành viên chậm tiến độ, PR chờ review lâu & Trung bình & Trung bình & Bảng phân công nội bộ có hạn từng PR; mỗi PR có người review cố định \\
RK6 & Bị đánh giá dùng lại sản phẩm của Assignment 1 & Thấp & Cao & Repository mới, lịch sử commit đi theo đúng thứ tự Task 1 đến Task 5 \\
\end{longtable}
}

Ở cuối dự án, rủi ro RK1 và RK2 đã được xử lý đúng theo cách giảm thiểu đã đề ra:
\begin{itemize}
\item MVP và pipeline CD có trước, chức năng bổ sung làm sau.
\item Production được dựng từ \code{v0.9.0}, trước khi giao diện khách hoàn chỉnh.
\end{itemize}

Rủi ro RK3 được kiểm chứng bằng các test đồng thời (mục \ref{ch:kiemthu}). Rủi ro RK5 vẫn xảy ra một phần: giao diện xong chậm hơn API, như Retrospective Sprint 2 đã ghi.

\subsection{Task 1.2 -- Yêu cầu chức năng và phi chức năng}
SRS gồm:
\begin{itemize}
\item 23 yêu cầu chức năng: 12 Must, 6 Should, 3 Could, 2 Won't.
\item 11 yêu cầu phi chức năng.
\end{itemize}

Mỗi yêu cầu có mã, mức ưu tiên MoSCoW và tiêu chí chấp nhận kiểm chứng được, phần lớn ở dạng mã HTTP, trạng thái dữ liệu hoặc số liệu đo được. Tiêu chí chấp nhận là căn cứ để viết test ở Task 4, 5 (xem ma trận truy vết, mục \ref{ch:truyvet}).

{\small
\begin{longtable}{@{}L{1.2cm}L{1.6cm}L{4.9cm}L{1.3cm}L{4.9cm}@{}}
\caption{Yêu cầu chức năng nhóm MVP (Must)}\label{tab:frmvp}\\
\toprule
\textbf{Mã} & \textbf{Nhóm} & \textbf{Yêu cầu} & \textbf{Ưu tiên} & \textbf{Tiêu chí chấp nhận} \\ \midrule \endfirsthead
\toprule \textbf{Mã} & \textbf{Nhóm} & \textbf{Yêu cầu} & \textbf{Ưu tiên} & \textbf{Tiêu chí chấp nhận} \\ \midrule \endhead
\bottomrule \endfoot
FR-01 & Menu & Hiển thị danh sách sản phẩm có tên, giá, ảnh và tình trạng còn hoặc hết & Must & \code{GET /products} trả danh sách; trang menu hiển thị đủ sản phẩm trong seed \\
FR-02 & Menu & Hiển thị trạng thái đang tải, danh sách rỗng và lỗi của menu & Must & Khi API chậm, trả rỗng hoặc trả lỗi, giao diện hiện đúng từng trạng thái \\
FR-03 & Chi tiết & Xem chi tiết sản phẩm; chọn size S/M/L, một hoặc nhiều topping (tối đa 5 mỗi ly), mức đường và đá (0, 30, 50, 70, 100\%) & Must & \code{GET /\allowbreak{}products/\allowbreak{}:id} trả đủ tùy chọn; giá cập nhật theo lựa chọn (S +0, M +5.000, L +10.000; cộng giá từng topping); không chọn được quá 5 topping hoặc topping đã hết nguyên liệu \\
FR-04 & Giỏ hàng & Thêm món, sửa số lượng, xóa món; hiển thị tổng tiền; tối đa 50 ly mỗi dòng & Must & Tổng tiền = $\Sigma$ (giá món + phụ phí size + giá topping) $\times$ số lượng; không tăng được quá 50 ly hoặc quá số ly nguyên liệu còn cho phép \\
FR-05 & Xác thực & Đăng ký, đăng nhập, phát JWT; mật khẩu băm bằng bcrypt & Must & Đăng nhập đúng nhận JWT; sai mật khẩu trả 401; database không lưu mật khẩu thô \\
FR-06 & Xác thực & Chặn đặt đơn khi chưa đăng nhập & Must & \code{POST /orders} không có JWT trả 401; giao diện chuyển sang đăng nhập và giữ giỏ hàng \\
FR-07 & Đặt đơn & Tạo đơn từ giỏ hàng ở trạng thái PENDING sau khi kiểm tra đủ nguyên liệu cho cả đơn & Must & Đơn tạo đúng món và giá; nếu thiếu nguyên liệu thì trả 400, báo món thiếu và không tạo đơn; tạo đơn chưa trừ nguyên liệu \\
FR-08 & Thanh toán & Thanh toán bằng Ví hoặc Thẻ mô phỏng. Thành công chuyển đơn sang PAID, thất bại chuyển sang PAYMENT\_FAILED; thất bại thì cho thử lại & Must & \code{POST /\allowbreak{}payments} đổi trạng thái đúng; thanh toán thất bại vẫn giữ giỏ hàng \\
FR-09 & Thanh toán & Gửi lại yêu cầu thanh toán với cùng Idempotency-Key không tạo thêm thanh toán & Must & Hai request cùng khóa nhận cùng kết quả; chỉ có một bản ghi \code{Payment} \\
FR-10 & Trạng thái đơn & Máy trạng thái đơn chặn mọi chuyển trạng thái sai & Must & Cho phép: PENDING $\rightarrow$ PAID, PAYMENT\_FAILED hoặc CANCELLED; PAYMENT\_FAILED $\rightarrow$ PAID hoặc CANCELLED; PAID $\rightarrow$ PREPARING hoặc CANCELLED; PREPARING $\rightarrow$ READY $\rightarrow$ COMPLETED. Chuyển sai trả 400 và đơn giữ nguyên \\
FR-11 & Lịch sử & Khách xem lịch sử và chi tiết đơn của mình & Must & \code{GET /\allowbreak{}orders/\allowbreak{}me} chỉ trả đơn của người đang đăng nhập; \code{GET /\allowbreak{}orders/\allowbreak{}:id} chỉ chủ đơn hoặc STAFF xem được, người khác nhận 404 \\
FR-12 & Tồn kho & Trừ nguyên liệu khi thanh toán thành công bằng cập nhật có điều kiện trong cùng transaction. Hủy đơn chỉ hoàn phần đã trừ & Must & Test N thanh toán đồng thời: số đơn PAID không vượt quá lượng nguyên liệu, tồn kho không âm. Nếu kho vừa hết đúng lúc thanh toán thì không trừ tiền, không trừ kho, đơn chuyển PAYMENT\_FAILED kèm lý do \\
\end{longtable}
}
{\small
\begin{longtable}{@{}L{1.2cm}L{1.6cm}L{4.9cm}L{1.3cm}L{4.9cm}@{}}
\caption{Yêu cầu chức năng bổ sung (Should, Could, Won't)}\label{tab:frbosung}\\
\toprule
\textbf{Mã} & \textbf{Nhóm} & \textbf{Yêu cầu} & \textbf{Ưu tiên} & \textbf{Tiêu chí chấp nhận} \\ \midrule \endfirsthead
\toprule \textbf{Mã} & \textbf{Nhóm} & \textbf{Yêu cầu} & \textbf{Ưu tiên} & \textbf{Tiêu chí chấp nhận} \\ \midrule \endhead
\bottomrule \endfoot
FR-13 & Khuyến mãi & Áp mã khuyến mãi hợp lệ để giảm giá & Should & Mã hợp lệ giảm đúng số tiền; mã sai, hết hạn hoặc hết lượt bị từ chối kèm lý do; giá cuối do máy chủ tính \\
FR-14 & Điểm thưởng & Cộng điểm khi đơn chuyển PAID; trừ lại khi hủy đơn đã PAID; mỗi lần thay đổi điểm ghi một giao dịch & Should & \code{User.points} tăng 1 điểm cho mỗi 10.000đ và giảm lại khi hủy; có bản ghi \code{LoyaltyTransaction} tương ứng (EARN, REVERSAL) \\
FR-15 & Nhân viên & Màn hình nhân viên xem đơn theo trạng thái và chuyển trạng thái & Should & \code{PATCH /\allowbreak{}orders/\allowbreak{}:id/\allowbreak{}status} chỉ STAFF gọi được (CUSTOMER nhận 403); trang nhân viên chỉ STAFF vào được \\
FR-16 & Nguyên liệu & Định lượng nguyên liệu theo công thức: công thức size S, size M và L nhân 125\% và 150\%, topping có định lượng riêng, đường và đá theo mức chọn & Should & Lượng nguyên liệu của một ly khớp công thức và hệ số (làm tròn lên); món hết nguyên liệu hiện "Hết" và không thêm vào giỏ được; trang chi tiết hiện "có thể đặt thêm tối đa N ly" mà không lộ số tồn kho \\
FR-17 & Nguyên liệu & Nhân viên xem tồn kho nguyên liệu, cờ "Sắp hết" và "Hết", và nhập thêm (restock) & Should & \code{GET /\allowbreak{}ingredients} và \code{PATCH /\allowbreak{}ingredients/\allowbreak{}:id/\allowbreak{}restock} chỉ STAFF gọi được; số lượng nhập phải dương \\
FR-18 & Thanh toán & Thanh toán bằng QR qua cổng mô phỏng: khách quét mã bằng điện thoại, chọn tài khoản thử nghiệm và nhập PIN & Should & Xác nhận đúng PIN và đủ số dư thì đơn chuyển PAID và trang đặt hàng tự chuyển sang xác nhận. Sai PIN, thiếu số dư hoặc QR hết hạn thì không trừ tiền. Xác nhận lặp lại chỉ thành công một lần \\
FR-19 & Đổi quà & Khách đổi điểm lấy quà và nhận mã quà để dùng tại quầy & Could & Đủ điểm thì trừ điểm và trả mã quà; không đủ điểm thì từ chối; đổi đồng thời không làm âm điểm \\
FR-20 & Đặt đơn & Tự hủy đơn chưa thanh toán sau thời gian cấu hình được (mặc định 15 phút) & Could & Đơn PENDING quá hạn chuyển CANCELLED và trả lại lượt dùng mã khuyến mãi; không đụng tới kho \\
FR-21 & Thanh toán & Thanh toán qua VNPay sandbox (môi trường thử nghiệm, không dùng tiền thật) & Could & Thanh toán thành công trên sandbox thì đơn chuyển PAID qua Return URL; máy chủ kiểm tra chữ ký HMAC và số tiền trước khi chốt đơn; sai chữ ký hoặc sai số tiền thì từ chối và đơn giữ nguyên trạng thái; thông tin cấu hình sandbox chỉ nằm trong biến môi trường \\
FR-22 & Quản trị & Giao diện quản lý sản phẩm, khuyến mãi, báo cáo doanh thu & Won't & Không làm trong phiên bản này \\
FR-23 & Thanh toán & Kết nối cổng thanh toán thật (tiền thật) và cổng MoMo & Won't & Không làm trong phiên bản này \\
\end{longtable}
}

Các yêu cầu phi chức năng tập trung vào năm nhóm:
\begin{itemize}
\item \textbf{Hiệu năng:} NFR-01.
\item \textbf{Bảo mật:} NFR-02 đến NFR-04.
\item \textbf{Khả dụng:} NFR-05, NFR-06.
\item \textbf{Triển khai, vận hành:} NFR-07 đến NFR-09.
\item \textbf{Độ tin cậy, bảo trì:} NFR-10, NFR-11.
\end{itemize}

NFR-08 yêu cầu có ít nhất một lần pipeline phát hiện lỗi. Yêu cầu này được kiểm chứng bằng PR \#10 cố ý làm test lỗi (mục \ref{ch:sprint1}) và lần CD \code{v0.9.1} bị đỏ (mục \ref{ch:sprint2}).

{\small
\begin{longtable}{@{}L{1.3cm}L{1.6cm}L{4.8cm}L{1.3cm}L{4.9cm}@{}}
\caption{Yêu cầu phi chức năng}\label{tab:nfr}\\
\toprule
\textbf{Mã} & \textbf{Nhóm} & \textbf{Yêu cầu} & \textbf{Ưu tiên} & \textbf{Tiêu chí chấp nhận} \\ \midrule \endfirsthead
\toprule \textbf{Mã} & \textbf{Nhóm} & \textbf{Yêu cầu} & \textbf{Ưu tiên} & \textbf{Tiêu chí chấp nhận} \\ \midrule \endhead
\bottomrule \endfoot
NFR-01 & Hiệu năng & \code{GET /products} phản hồi nhanh với dữ liệu seed & Should & 95\% trong 20 yêu cầu liên tiếp dưới 500 ms khi chạy cục bộ bằng Docker Compose \\
NFR-02 & Bảo mật & Mật khẩu và PIN được băm; JWT có hạn; kiểm tra phân quyền theo vai trò & Must & Có test xác thực và phân quyền; không có mật khẩu hoặc PIN thô trong database hay log \\
NFR-03 & Bảo mật & Secret không nằm trong repository hoặc log CI & Must & \code{.env} có trong \code{.gitignore}; secret của production lưu trên nền tảng triển khai hoặc GitHub Secrets \\
NFR-04 & Bảo mật & Kiểm tra dữ liệu đầu vào ở cả frontend và backend & Must & Dữ liệu sai hoặc trường lạ trả 400 với thông báo rõ ràng \\
NFR-05 & Khả dụng & Giao diện mobile-first & Must & Dùng được trên viewport 360--414 px, không cuộn ngang \\
NFR-06 & Khả dụng & Thông báo lỗi bằng tiếng Việt, dễ hiểu, không lộ chi tiết kỹ thuật & Should & Không hiển thị stack trace hay thông báo lỗi gốc của máy chủ \\
NFR-07 & Triển khai & Chạy toàn bộ hệ thống bằng \code{docker compose up} & Must & Người ngoài nhóm chạy được theo README \\
NFR-08 & Triển khai & CI chạy lint, test và build cho mọi PR; CD build image, triển khai và smoke test, có cách rollback & Must & Có log pipeline chạy thành công và ít nhất một lần pipeline phát hiện lỗi \\
NFR-09 & Vận hành & Có endpoint \code{GET /health} và log & Must & \code{/health} trả \code{status: ok}; smoke test gọi \code{/health} sau mỗi lần deploy \\
NFR-10 & Độ tin cậy & Migration và seed tái tạo được database & Must & Chạy migration và seed trên database trống luôn cho cùng kết quả \\
NFR-11 & Bảo trì & Mã nguồn có lint, định dạng thống nhất và test tự động & Should & CI chạy lint và test; có unit test và integration test cho nghiệp vụ lõi \\
\end{longtable}
}

\subsubsection*{Lịch sử thay đổi yêu cầu}
Ngày 17/11/2026, FR-03 được ghi rõ ``chọn một hoặc nhiều topping, tối đa 5 mỗi ly''. Lý do: giao diện \code{v1.0.0} chỉ cho chọn một topping, trong khi API và nghiệp vụ cho phép nhiều. Thay đổi được ghi trong mục ``Lịch sử thay đổi'' của \code{srs.md} kèm số PR (\#56).

\subsection{Use case tổng quát}
Hệ thống có ba actor:
\begin{itemize}
\item \textbf{Khách hàng:} xem menu, xem chi tiết, quản lý giỏ hàng, đăng ký, đăng nhập, đặt đơn, thanh toán, xem lịch sử, đổi quà.
\item \textbf{Nhân viên:} xử lý đơn (cập nhật trạng thái), hủy đơn, quản lý nguyên liệu.
\item \textbf{Cổng thanh toán (mô phỏng):} tham gia use case thanh toán.
\end{itemize}

\hinh{usecase-overview.pdf}{0.95\linewidth}{Sơ đồ use case tổng quát (\code{docs/\allowbreak{}requirements/\allowbreak{}usecase-overview.\allowbreak{}puml})}{fig:ucoverview}

\subsection{Task 1.3 -- Đặc tả use case trọng tâm UC-01}
Nhóm chọn \textbf{UC-01 ``Đặt và thanh toán đơn hàng''} làm use case trọng tâm. Đây là luồng tạo ra doanh thu. Sai ở luồng này gây mất tiền, trừ tiền hai lần hoặc bán vượt nguyên liệu, nên mọi cơ chế khó của hệ thống (transaction, idempotency, máy trạng thái) đều nằm trên luồng này.

\hinh{usecase-checkout.pdf}{0.9\linewidth}{Sơ đồ use case cho phạm vi UC-01 (\code{docs/\allowbreak{}requirements/\allowbreak{}usecase-checkout.\allowbreak{}puml})}{fig:uccheckout}

{\small
\begin{longtable}{@{}L{3.2cm}L{10.7cm}@{}}
\caption{Thông tin chung của UC-01}\label{tab:ucinfo}\\
\toprule
\textbf{Mục} & \textbf{Nội dung} \\ \midrule \endfirsthead
\toprule \textbf{Mục} & \textbf{Nội dung} \\ \midrule \endhead
\bottomrule \endfoot
Mã & UC-01 \\
Tên & Đặt và thanh toán đơn hàng \\
Actor chính & Khách hàng (CUSTOMER) \\
Actor phụ & Cổng thanh toán (mô phỏng) \\
Trigger & Khách bấm "Đặt đơn" trong giỏ hàng \\
Yêu cầu liên quan & FR-04, FR-05, FR-06, FR-07, FR-08, FR-09, FR-10, FR-12, FR-13, FR-14, FR-16, FR-18, NFR-02, NFR-04 \\
\end{longtable}
}

\textbf{Tiền điều kiện:}
\begin{itemize}
\item Khách đã đăng nhập và có JWT hợp lệ.
\item Giỏ hàng có ít nhất một món.
\end{itemize}

\textbf{Hậu điều kiện:}
\begin{itemize}
\item \textbf{Thành công:}
  \begin{itemize}
  \item Đơn ở trạng thái PAID; nguyên liệu đã bị trừ; có bản ghi \code{Payment}; điểm thưởng đã được cộng.
  \item Khách thấy trang xác nhận đơn.
  \end{itemize}
\item \textbf{Thanh toán thất bại:} đơn ở PAYMENT\_FAILED, giỏ hàng được giữ; khách có thể thử lại. Nhân viên có quyền hủy đơn theo phân quyền.
\item \textbf{Thiếu nguyên liệu khi đặt:} không tạo đơn, tồn kho không đổi.
\end{itemize}

\textbf{Luồng chính:}
\begin{enumerate}
\item Khách mở giỏ hàng và bấm ``Đặt đơn''.
\item Giao diện gửi \code{POST /orders} kèm JWT.
\item API kiểm tra JWT và dữ liệu đầu vào.
\item API mở một transaction và thực hiện:
  \begin{itemize}
  \item cộng dồn nhu cầu nguyên liệu của cả đơn theo công thức, size, topping, mức đường và đá;
  \item kiểm tra đủ nguyên liệu và ghi lại nhu cầu (chưa trừ kho);
  \item tạo \code{Order} (PENDING) và các \code{OrderItem}, áp khuyến mãi nếu có;
  \item commit.
  \end{itemize}
\item API trả mã đơn; giao diện chuyển sang màn hình thanh toán.
\item Khách chọn Ví hoặc Thẻ và xác nhận.
\item Giao diện gửi \code{POST /\allowbreak{}payments} kèm \code{Idempotency-Key}.
\item API thực hiện transaction thứ nhất: trừ nguyên liệu bằng cập nhật có điều kiện (\code{stock >= lượng cần}); nếu đủ thì commit, nếu thiếu thì rollback.
\item API gọi cổng thanh toán mô phỏng. Nếu cổng trả thành công, API mở transaction thứ hai để:
  \begin{itemize}
  \item ghi \code{Payment};
  \item chuyển đơn sang PAID;
  \item cộng điểm thưởng;
  \item commit.
  \end{itemize}
\item API trả kết quả; giao diện hiện trang xác nhận đơn.
\end{enumerate}

{\small
\begin{longtable}{@{}L{0.9cm}L{4.4cm}L{8.6cm}@{}}
\caption{Luồng thay thế của UC-01}\label{tab:ucalt}\\
\toprule
\textbf{\#} & \textbf{Điều kiện} & \textbf{Xử lý} \\ \midrule \endfirsthead
\toprule \textbf{\#} & \textbf{Điều kiện} & \textbf{Xử lý} \\ \midrule \endhead
\bottomrule \endfoot
A1 & Khách nhập mã khuyến mãi hợp lệ trước bước 1 & Giá được giảm; tổng tiền cập nhật trước khi đặt đơn \\
A2 & Khách gửi lại yêu cầu thanh toán với cùng Idempotency-Key, ví dụ khi mạng chập chờn & API trả lại kết quả cũ, không tạo thanh toán thứ hai \\
A3 & Ở bước 6 khách chọn QR & Thay cho bước 6--7:  1. API tạo mã QR có hạn dùng.  2. Khách quét mã bằng điện thoại, chọn tài khoản thử nghiệm và nhập PIN.  3. API thực hiện bước 8; cổng mô phỏng kiểm tra PIN và số dư thay cho lời gọi cổng ở bước 9.  4. Trang đặt hàng hỏi trạng thái định kỳ và tự chuyển sang trang xác nhận.  Phần ghi nhận của bước 9 và bước 10 giữ nguyên \\
A4 & STAFF hủy đơn ở trạng thái cho phép (PENDING, PAYMENT\_FAILED hoặc PAID) & Đơn chuyển CANCELLED. Chỉ phần nguyên liệu đã trừ được hoàn lại; đơn chưa thanh toán không đụng tới kho. Điểm đã cộng bị trừ lại \\
\end{longtable}
}
{\small
\begin{longtable}{@{}L{0.9cm}L{4.4cm}L{8.6cm}@{}}
\caption{Luồng ngoại lệ của UC-01}\label{tab:ucex}\\
\toprule
\textbf{\#} & \textbf{Điều kiện} & \textbf{Xử lý} \\ \midrule \endfirsthead
\toprule \textbf{\#} & \textbf{Điều kiện} & \textbf{Xử lý} \\ \midrule \endhead
\bottomrule \endfoot
E1 & Chưa đăng nhập hoặc JWT hết hạn (bước 2--3) & API trả 401; giao diện chuyển sang đăng nhập rồi quay lại giỏ hàng; giỏ hàng được giữ \\
E2 & Dữ liệu không hợp lệ (bước 3) & API trả 400 với thông báo rõ ràng \\
E3 & Không đủ nguyên liệu khi tạo đơn (bước 4) & Transaction rollback; không tạo đơn; báo món nào không đủ nguyên liệu \\
E4 & Kho vừa hết do đơn khác thanh toán trước (bước 8) & Transaction bị hủy trước khi gọi cổng thanh toán: không trừ tiền, không trừ kho; đơn chuyển PAYMENT\_FAILED kèm lý do; khách thử lại hoặc hủy đơn \\
E5 & Mã khuyến mãi sai, hết hạn hoặc hết lượt & Từ chối mã kèm lý do; khách vẫn đặt được đơn không dùng mã \\
E6 & Cổng thanh toán trả thất bại (bước 9), hoặc QR sai PIN, thiếu số dư hay hết hạn (A3) & Transaction rollback nên kho không đổi; không trừ tiền; Ví và Thẻ thì đơn chuyển PAYMENT\_FAILED. Giao diện báo lỗi, giữ giỏ hàng và cho thử lại \\
E7 & Yêu cầu chuyển trạng thái không hợp lệ & API trả 400; đơn giữ nguyên trạng thái \\
\end{longtable}
}

Luồng ngoại lệ E4 là điểm quan trọng nhất của thiết kế. Hai khách cùng thanh toán khi kho chỉ còn đủ cho một đơn. Kho bị trừ bằng cập nhật có điều kiện \emph{trước} khi gọi cổng thanh toán. Vì vậy người đến sau không bị trừ tiền, đơn chuyển PAYMENT\_FAILED kèm lý do. Cơ chế này được ghi ở ADR 005 và kiểm chứng bằng integration test INT-P06, INT-P07.
